\chapter{Background}
\label{chap:background}

\section{Mobile Application Platform}

The prototype is built as a React Native application. React Native enables native mobile interfaces using JavaScript and React concepts \cite{react-native}. Expo, an open-source framework built on top of React Native, was selected because the project needs camera access, image import, SQLite persistence, secure storage, and web support without introducing native modules or custom development clients. Expo Router provides file-based routing for React Native and web applications, which matches the required Capture, Confirm and Edit, History, and Settings screens \cite{expo-router}.

\section{Vision-Language Extraction}

The core recognition problem is not only optical character recognition. Product labels contain layout, abbreviations, model numbers, manufacturer names, country information, barcodes, and sometimes technical identifiers. A pure OCR pipeline can transcribe characters, but would still need a downstream mapping layer to decide which recognized text belongs to which registration field. The implementation therefore uses a vision-language extraction step to map visible label evidence directly into a structured JSON schema.

This distinction matters because transcription and registration are different tasks. OCR can make label text machine-readable, but inventory registration also requires field selection, normalization, and omission of values that are not visible or reliable. The project therefore treats model output as draft structured data, while deterministic barcode decoding, schema validation, and human review remain separate control points.

Planning still considered OCR and text-recognition approaches. CRNN represents neural scene-text recognition for image-based sequences \cite{crnn}, while Tesseract represents a classical OCR engine that was relevant as a baseline OCR option \cite{tesseract}. Gemini document processing and ML Kit Text Recognition v2 were also relevant practical options during planning \cite{gemini-doc-processing,mlkit-text-recognition}. These sources informed the design direction, but the implementation keeps deterministic barcode decoding separate and uses a VLM-first extraction boundary for schema mapping.

\section{Architecture Documentation}

The project uses C4-style architecture thinking, where the system is described through increasingly detailed views such as context, containers, components, and code \cite{c4-model}. The diagrams in this report are generated from PlantUML sources and rendered for inclusion in the report. Only the rendered diagram assets are included in the appendix and final report. The PlantUML source files remain project working files.

\section{Related Implementation Constraints}

The Expo managed workflow shaped the implementation. It supports portability across iOS, Android, and web, but it also discourages custom native model runtimes. For that reason, local vision-language inference remains future work. The current prototype prioritizes a reliable remote provider boundary, local persistence, human review, and replay-safe submission.